\documentclass[openany]{book}
\makeatletter
\renewcommand{\@chapapp}{Lecture}
\makeatother

\usepackage{hyperref}
\hypersetup{colorlinks}

\usepackage{listings}
\usepackage[usenames, dvipsnames]{color}

\definecolor{mygreen}{rgb}{0,0.6,0}
\definecolor{mygray}{rgb}{0.5,0.5,0.5}
\definecolor{mymauve}{rgb}{0.58,0,0.82}

\lstset{ %
  backgroundcolor=\color{white},   % choose the background color; you must add \usepackage{color} or \usepackage{xcolor}; should come as last argument
  basicstyle=\footnotesize\ttfamily,        % the size of the fonts that are used for the code
  breakatwhitespace=false,         % sets if automatic breaks should only happen at whitespace
  breaklines=true,                 % sets automatic line breaking
  captionpos=b,                    % sets the caption-position to bottom
  commentstyle=\color{mygreen},    % comment style
  deletekeywords={...},            % if you want to delete keywords from the given language
  escapeinside={\%*}{*)},          % if you want to add LaTeX within your code
  extendedchars=true,              % lets you use non-ASCII characters; for 8-bits encodings only, does not work with UTF-8
  frame=single,	                   % adds a frame around the code
  keepspaces=true,                 % keeps spaces in text, useful for keeping indentation of code (possibly needs columns=flexible)
  keywordstyle=\color{blue},        % keyword style
%  language=Octave,                 % the language of the code
  language=Scilab,
  morekeywords={\$, \%F, \%T, \%e,
  \%eps, \%f, \%fftw, \%gui,
  \%i, \%inf, \%io, \%nan,
  \%pi, \%s, \%t, \%tk,
  surf, scf, clf, meshgrid, }, %keywordstyle=\color{BurntOrange},classoffset=0,           % if you want to add more keywords to the set
  numbers=left,                    % where to put the line-numbers; possible values are (none, left, right)
  numbersep=5pt,                   % how far the line-numbers are from the code
  numberstyle=\tiny\color{mygray}, % the style that is used for the line-numbers
  rulecolor=\color{black},         % if not set, the frame-color may be changed on line-breaks within not-black text (e.g. comments (green here))
  showspaces=false,                % show spaces everywhere adding particular underscores; it overrides 'showstringspaces'
  showstringspaces=false,          % underline spaces within strings only
  showtabs=false,                  % show tabs within strings adding particular underscores
  stepnumber=2,                    % the step between two line-numbers. If it's 1, each line will be numbered
  stringstyle=\color{mymauve},     % string literal style
  tabsize=2,	                   % sets default tabsize to 2 spaces
  title=\lstname                   % show the filename of files included with \lstinputlisting; also try caption instead of title
}
\begin{document}
\frontmatter
% Title page
\begin{titlepage}
\newcommand{\HRule}{\rule{\linewidth}{0.5mm}} % Defines a new command for the horizontal lines, change thickness here

\center % Center everything on the page

%----------------------------------------------------------------------------------------
%	HEADING SECTIONS
%----------------------------------------------------------------------------------------

\textsc{\LARGE City University of Hong Kong}\\[1.5cm] % Name of your university/college
\textsc{\Large AP Courses Review Notes}\\[0.5cm] 		% Major heading such as course name
\textsc{\large {AP3114}}\\[0.5cm] 						% Minor heading such as course title

%----------------------------------------------------------------------------------------
%	TITLE SECTION
%----------------------------------------------------------------------------------------

\HRule \\[0.4cm]
{ \huge \bfseries  \textsf{Computational Methods for Physicists and Materials Engineers} }\\[0.4cm] % Title of your document
\rightline{{\bfseries \textsf{Version 1.0}} }
\HRule \\[1.5cm]

%----------------------------------------------------------------------------------------
%	AUTHOR SECTION
%----------------------------------------------------------------------------------------



% If you don't want a supervisor, uncomment the two lines below and remove the section above
%\Large \emph{Author:}\\
%John \textsc{Smith}\\[3cm] % Your name

%----------------------------------------------------------------------------------------
%	DATE SECTION
%----------------------------------------------------------------------------------------


{\large \today}\\[2cm] % Date, change the \today to a set date if you want to be precise

%----------------------------------------------------------------------------------------
%	LOGO SECTION
%----------------------------------------------------------------------------------------
\vspace{5cm}
%\includegraphics[scale=0.5]{figures/CityU_full_logo_2015.pdf}\\[1cm] % Include a department/university logo - this will require the graphicx package
%\includegraphics[scale=0.5]{figures/CityU_logo_2015.eps}\\[1cm] % Include a department/university logo - this will require the graphicx package


%----------------------------------------------------------------------------------------

\vfill % Fill the rest of the page with whitespace
\end{titlepage}

\noindent

This document is free; you can redistribute it and/or modify it under the terms of 
the GNU General Public License as published by the Free Software Foundation; either 
version 2 of the License, or (at your option) any later version.

This document is distributed in the hope that it will be useful, but \emph{without any 
warranty}; without even the implied warranty of \emph{merchantability} or \emph{fitness 
for a particular purpose}. See the GNU General Public License for more details.

You should have received a copy of the GNU General Public License along with this document; 
if not, write to the Free Software Foundation, Inc., 675 Mass Ave, Cambridge, MA 02139, USA.

All \LaTeX{} (\texttt{.tex}) files of this document can be accessed from
\url{https://github.com/zzw42/review-notes-cityu}.


\tableofcontents

\chapter*{Preface}
\addcontentsline{toc}{chapter}{Preface}
The first few exercises are from \emph{Introduction to Scilab} by Micha\"{e}l Baudin.
The homepage of this document is 
\url{http://forge.scilab.org/index.php/p/docintrotoscilab/}, where the source can also be
found.

The note on cubic equation by Urs Oswald is also 
helpful. \url{http://www.ursoswald.ch/download/CUBIC.pdf}.

These {\tt .sci} files have been tested under Scilab 5.52.
\mainmatter
\chapter{Typical Problems in Physics and Materials Engineering
\& Introduction to Scilab}
\section{Homework}
\lstinputlisting{./personal_answer_Homework/lecture1_exe.sce}


\chapter{Matrix}
\section{Lecture}
\lstinputlisting{./lecture/lecture_2_matrix.sce}
\section{Homework}
\lstinputlisting{./personal_answer_Homework/lecture2_exe.sce}


\chapter{Basic Programming}
\section{Lecture}
\lstinputlisting{./lecture/lecture_3_basic_programming.sce}
\section{Homework}
\lstinputlisting{./personal_answer_Homework/lecture3_exe.sce}


\chapter{Program II \& Plotting I}
\section{Lecture}
\lstinputlisting{./lecture/lecture_4_programmingii_plottingi.sce}
\section{Homework}
\lstinputlisting{./personal_answer_Homework/lecture4_exe.sce}


\chapter{Programming III \& Plotting II}
\section{Lecture}
\lstinputlisting{./lecture/lecture_5_programmingiii_plottingii.sce}
\section{Homework}
\lstinputlisting{./personal_answer_Homework/lecture5_exe1.sci}
\lstinputlisting{./personal_answer_Homework/lecture5_exe2.sci}
\lstinputlisting{./personal_answer_Homework/lecture5_exe3.sci}
\lstinputlisting{./personal_answer_Homework/lecture5_exe4.sci}
\lstinputlisting{./personal_answer_Homework/lecture5_exe5.sci}
\lstinputlisting{./personal_answer_Homework/lecture5_exe6.sci}


\chapter{Random Number}
\section{Lecture}
\lstinputlisting{./lecture/lecture_6_random-number.sce}
\section{Homework}
\lstinputlisting{./personal_answer_Homework/lecture6_exe1.sci}
\lstinputlisting{./personal_answer_Homework/lecture6_exe1b.sci}
\lstinputlisting{./personal_answer_Homework/lecture6_exe2.sci}
\lstinputlisting{./personal_answer_Homework/lecture6_exe3.sci}
\lstinputlisting{./personal_answer_Homework/lecture6_exe4.sci}
\lstinputlisting{./personal_answer_Homework/lecture6_exe5.sci}


\chapter{Statistical Description and Analysis of Data}
\section{Lecture}
\lstinputlisting{./lecture/lecture_7_statistics.sce}
\section{Homework}
\lstinputlisting{./personal_answer_Homework/lecture7_exe1.sci}
\lstinputlisting{./personal_answer_Homework/lecture7_exe2.sci}


\chapter{Solving Non-linear Equations}
\section{Lecture}
\lstinputlisting{./lecture/lecture_8_9_solving_nonlinear_equations.sce}

\section{Homework}
\paragraph{Exercise 1}
Calculate all the cubic roots of below numbers in complex field:
\begin{enumerate}
\item $1$
\item $i$
\item $5+3i$
\item $2 e ^{i\pi/2}$
\item $-2-3i$
\end{enumerate}
\lstinputlisting{./personal_answer_Homework/lecture8_exe1.sci}

\paragraph{Exercise 2}
Write a Scilab program {\tt quadroots} to compute and print the roots of a quadratic equation $a x^2+ b x + c = 0$ according to  the previous pseudo-code.
It should run with a command like
\begin{center}
{\tt quadroots(1,3,2)}
\end{center}
for the case of $x^2+ 3x + 2 = 0$.

Test {\tt quadroots} on the following examples, checking the results in each case.
\begin{enumerate}
\item $x^2+ 1 = 0$
\item $0x^2+ 2x + 1 = 0$
\item $x^2+ 3x + 2 = 0$
\item $4x^2+ 24x + 36 = 0$
\end{enumerate}
\lstinputlisting{./personal_answer_Homework/lecture8_exe2.sci}

\paragraph{Exercise 3}
Write a Scilab program {\tt cubicroots} to compute and print the roots of a quadratic equation $ax^3+ bx^2 + cx +d = 0$ according to  the previous pseudo-code.
It should run with a command like
\begin{center}
{\tt cubicroots(4,5,1,2)}
\end{center}
for the case of $4x^3+5x^2+ 1x + 2 = 0$.

Test {\tt cubicroots} on the following examples, checking the results in each case.
\lstinputlisting{./personal_answer_Homework/lecture8_exe3.sci}

\paragraph{Exercise 4}
Find the roots of below polynomial using the function {\tt roots}:
\lstinputlisting{./personal_answer_Homework/lecture8_exe4.sci}

\paragraph{Exercise 5}

\lstinputlisting{./personal_answer_Homework/lecture8_exe5.sci}


\setcounter{chapter}{9}

\chapter{Solving Differential Equations I}
\section{Lecture}
\lstinputlisting{./lecture/lecture_10_solving_differential_equations_i.sce}
\section{Homework}
\lstinputlisting{./personal_answer_Homework/lecture10_exe1.sci}
\lstinputlisting{./personal_answer_Homework/lecture10_exe2.sci}
\lstinputlisting{./personal_answer_Homework/lecture10_exe3.sci}
\lstinputlisting{./personal_answer_Homework/lecture10_exe4.sci}


\chapter{Solving Differential Equations II}
\section{Lecture}
\lstinputlisting{./lecture/lecture_11_solving_differential_equations_ii.sce}



\chapter{Fourier Series}
\section{Lecture}
\lstinputlisting{./lecture/lecture_12_fourier.sce}
\section{Homework}
\lstinputlisting{./personal_answer_Homework/lecture12_exe.sci}


\appendix
\chapter{Mid-Term Exam}
\lstinputlisting{personal_answer_midterm.sci}


\chapter{Final Exam}
\lstinputlisting{personal_answer_final.sce}
\end{document}
